%% 18장 언어의 경계
\chapter{The Language Boundary}
\label{ch:language-boundary}

\abstract{과학 계산 라이브러리의 대부분은 두 언어로 된다. 사용자는 Python으로 쓰고, 무거운 계산은 C나 C++로 된 확장 모듈이 한다. 이 장은 그 경계의 양쪽과 경계 자체를 원리 수준에서 다룬다. C++ 쪽에서는 자료형마다 코드를 한 번씩 찍어 내는 템플릿, 자원의 수명을 객체의 수명에 묶는 RAII, 그리고 컴파일러가 ``일어나지 않는다''고 가정하는 정의되지 않은 동작과 그 가운데 수치 코드에 특히 위험한 별칭 규칙을 본다. Python 쪽에서는 모든 것이 객체이고 참조 계수로 수명이 관리되는 CPython의 객체 모델과 전역 인터프리터 잠금을 본다. 마지막으로 두 세계 사이에서 메모리를 복사 없이 넘기는 버퍼 프로토콜과, 작은 호출에서 경계를 넘는 비용이 계산의 비용보다 커지는 문제를 다룬다.}

\section{The C++ Side}
\label{sec:cpp-side}

\subsection{Templates and Types}

배열 라이브러리의 원소별 연산은 자료형마다 다른 기계어가 필요하다. 정수 덧셈과 부동소수점 덧셈은 다른 명령이고, 반복문의 벡터화도 자료형에 따라 다르다. 그렇다고 연산마다 자료형마다 코드를 손으로 쓰면 같은 논리가 수십 번 되풀이된다. C++의 \term{템플릿}{template}은 자료형을 매개변수로 하는 코드를 한 번 쓰고, 컴파일러가 쓰이는 자료형마다 따로 \term{인스턴스화}{instantiation}해서 각자에 맞는 기계어를 만들게 한다\cite{vandevoorde2017templates}. 실행 중에만 알 수 있는 배열의 자료형은 \code{switch} 하나로 컴파일 시점의 자료형으로 바꾸어 알맞은 인스턴스를 고른다.

\begin{instardust}[\code{csrc/base/dtype.hpp}]
\code{visit\_dtype(d, f)}가 이 다리이다. 실행 중의 자료형 값 \code{d}를 보고 \code{f(std::type\_identity<T>\{\})}를 부르는데, \code{T}는 \code{bool}, \code{std::int64\_t}, \code{double}, \code{std::complex<double>} 가운데 하나이다. 그래서 하나의 템플릿 함수 몸체가 자료형마다 한 번씩 인스턴스화되고 실행 중에 고르게 된다. 원소별 연산은 연산마다 구조체 하나에 ``이 자료형의 반복문이 있는가''와 원소 하나의 함수를 적어 두고, 이것을 반복문으로 바꾸는 일은 템플릿이 맡는다.
\end{instardust}

\subsection{RAII}

메모리, 파일, 잠금 같은 자원은 얻은 뒤 반드시 돌려주어야 한다. 함수에 돌아가는 길이 여럿이고 예외가 날 수도 있으면 돌려주는 코드를 빠뜨리기 쉽다. C++의 관용구 \term{RAII}{resource acquisition is initialization}는 자원을 객체의 생성자에서 얻고 소멸자에서 돌려준다\cite{stroustrup2013cpp}. 객체는 범위를 벗어날 때 어떤 길로 벗어나든 소멸자가 불리므로, 자원의 수명이 객체의 수명에 묶여 빠짐없이 돌려진다. 참조 계수로 공유되는 메모리를 관리하는 \code{std::shared\_ptr}도 RAII이다. 마지막 소유자가 사라질 때 메모리가 풀린다.

\begin{instardust}[\code{csrc/parallel/pool.hpp}]
14장에서 본 GIL 풀기가 RAII의 예이다. \code{GilRelease} 객체를 만들면 그 생성자가 GIL을 풀고, 객체가 범위를 벗어나면 소멸자가 다시 잡는다. 그 사이에서 예외가 나도 GIL은 반드시 다시 잡힌다. 배열의 메모리인 \code{Buffer}는 \code{shared\_ptr}로 공유되고, 버퍼 프로토콜로 바깥에서 가져온 메모리는 소멸자에서 \code{Py\_buffer}를 돌려준다. 마지막 소유자가 GIL이 없는 일꾼 스레드에서 사라질 수 있으므로 그 소멸자는 GIL을 잡은 뒤 돌려준다고 설계 문서는 적는다.
\end{instardust}

\subsection{Undefined Behavior}

C++ 표준은 몇 가지 상황에서 프로그램의 동작을 아예 정의하지 않는다. 부호 있는 정수의 넘침(2장), 범위를 벗어난 부동소수점의 정수 변환(8장), 데이터 경쟁(14장), 배열 범위 밖의 접근, 그리고 아래에서 볼 별칭 규칙의 위반이 그것이다. ``정의되지 않음''은 ``기계가 하는 대로''가 아니라 ``컴파일러가 그런 일이 일어나지 않는다고 가정해도 된다''는 뜻이다\cite{iso14882-2020}. 그래서 컴파일러는 그 가정에서 나오는 최적화를 하고, 가정이 깨지는 입력에서는 결과를 예측할 수 없다. 예컨대 \code{x + 1 > x}는 부호 있는 정수에서 넘침이 없다는 가정에 따라 언제나 참으로 바뀌어 사라질 수 있다\cite{lattner2011ub}. 수치 코드의 정의되지 않은 동작은 대부분 경계의 값, 곧 넘침, 무한대, NaN, 빈 배열에서 터지므로, 그런 값을 명시적으로 다루거나 정의된 동작이 보장되는 수단(부호 없는 정수, 포화 변환 명령)을 쓴다.

\subsection{Aliasing Rules}

두 포인터가 같은 메모리를 가리키는 것을 \term{별칭}{aliasing}이라 한다. 컴파일러는 별칭이 있을 수 있으면 한 포인터로 쓴 뒤 다른 포인터로 읽는 값을 다시 메모리에서 읽어야 하므로, 반복문의 값을 레지스터에 두거나 벡터화하기 어렵다. C++의 \term{엄격한 별칭 규칙}{strict aliasing rule}은 이 문제를 부분적으로 풀어 준다. 서로 다른 자료형의 포인터, 예컨대 \code{double*}과 \code{std::uint64\_t*}는 같은 객체를 가리키지 않는다고 가정해도 되는 것이다. 그 대가로, 부동소수점 수의 비트를 정수로 읽으려고 \code{reinterpret\_cast<std::uint64\_t*>(\&x)}처럼 쓰면 정의되지 않은 동작이 된다. 8장과 11장에서 본 비트 수준의 기교들은 모두 이런 읽기를 필요로 하므로, 표준이 허락하는 방법, 곧 \code{std::memcpy}로 바이트를 옮기거나 C++20의 \code{std::bit\_cast}를 쓴다. 두 방법 모두 최적화 뒤에는 레지스터 사이의 이동 하나로 끝나서 비용이 없다.

같은 자료형의 포인터 사이의 별칭은 컴파일러가 배제할 수 없다. 입력과 출력이 겹치지 않는다는 것을 프로그래머가 알 때는 \code{\_\_restrict} 같은 확장으로 컴파일러에 알려 벡터화를 돕는다. 17장에서 본 메모리 겹침의 판정이 이 지점과 맞닿아 있다. 겹치지 않는 것이 확인된 입력과 출력만 겹침 없는 커널에 넘긴다.

\section{The Python Side}
\label{sec:python-side}

\subsection{Objects and Reference Counting}

CPython에서는 정수 하나까지 모든 값이 힙에 할당된 객체이고, 객체마다 머리에 \term{참조 계수}{reference count}와 자료형을 가리키는 포인터가 있다\cite{python-c-api}. 어떤 변수나 컨테이너가 객체를 가리키게 되면 계수가 늘고, 그 참조가 사라지면 준다. 계수가 $0$이 되는 순간 객체는 즉시 해제된다. 순환 참조는 따로 도는 쓰레기 수집기가 찾는다. 참조 계수는 수명을 결정적으로 만들어 준다는 장점이 있지만, 참조를 주고받을 때마다 계수를 고쳐야 하므로 작은 객체를 많이 만드는 코드에서는 그 비용이 쌓인다. 확장 모듈은 객체를 받아 저장하거나 돌려줄 때 계수를 정확히 맞추어야 하고, 하나라도 틀리면 메모리가 새거나 이미 해제된 객체를 읽게 된다.

CPython은 $512$바이트 이하의 작은 객체를 위해 전용 할당기를 두어 시스템 \code{malloc}보다 훨씬 빨리 할당한다. 이 경계가 확장 모듈의 설계에 영향을 주기도 한다. 확장 객체가 $512$바이트를 넘으면 작은 객체 할당기를 쓰지 못하고 시스템 할당기로 가기 때문이다.

\subsection{The Global Interpreter Lock}

CPython의 인터프리터 상태와 참조 계수는 여러 스레드가 동시에 고치면 깨진다. 그래서 CPython은 한 번에 한 스레드만 Python 코드를 실행하게 하는 \term{전역 인터프리터 잠금}{global interpreter lock, GIL}을 둔다. GIL 덕분에 인터프리터의 내부는 단순해지지만, Python 스레드 여러 개가 동시에 계산할 수는 없다. 확장 모듈은 Python 객체를 건드리지 않는 긴 계산을 하는 동안 GIL을 풀어 다른 Python 스레드가 돌게 할 수 있으며, 14장에서 본 스레드 풀의 일꾼들은 GIL 없이 C++ 데이터만 다룬다. 최근에는 GIL을 선택적으로 없앤 CPython 빌드도 나왔는데\cite{pep703}, 그 환경에서는 확장 모듈이 스스로 스레드 안전성을 보장해야 한다.

\begin{instardust}[\code{csrc/parallel/pool.hpp}, \code{docs/design/frame.md} §5.3]
stardust의 규칙은 ``각 영역의 병렬 문턱값을 넘는 일은 스레드 수와 상관없이 GIL을 풀고 계산하고, 그 아래의 일은 GIL을 쥔 채 계산한다''이다. 작은 호출은 GIL을 풀고 잡는 비용조차 아깝기 때문이다. GIL을 푼 구간에서는 Python 객체를 건드릴 수 없으므로, 오류가 나면 14장의 \code{FirstRow}처럼 그 위치만 기록해 두었다가 GIL을 다시 잡은 뒤에 오류 메시지를 만든다. 데이터프레임은 GIL을 푼 계산과 다른 Python 스레드의 대입이 겹칠 수 있다는 점까지 고려한다. 진행 중인 연산이 옛 열 객체를 잡고 있으므로 대입은 17장의 Copy-on-Write 규칙에 따라 복사하게 되는데, 열의 버퍼를 그 자리에서 고치는 행 대입 하나만은 그 틈을 막으려고 크기와 상관없이 GIL을 잡은 채 돈다.
\end{instardust}

\section{Passing Data across the Boundary}
\label{sec:crossing}

\subsection{The Buffer Protocol}

Python 객체가 담은 메모리를 다른 확장 모듈이 복사 없이 읽으려면 그 메모리의 주소와 모양과 보폭과 자료형을 알려 주는 약속이 필요하다. 그것이 \term{버퍼 프로토콜}{buffer protocol}이다\cite{pep3118}. 메모리를 내보내는 쪽은 \code{Py\_buffer} 구조체에 주소, 차원, 모양, 보폭(17장), 그리고 원소의 자료형을 나타내는 형식 문자열을 채워 주고, 가져가는 쪽이 다 쓰면 돌려준다. 형식 문자열은 \code{struct} 모듈의 표기를 따라 \code{d}는 binary64, \code{q}는 $64$비트 정수, \code{?}는 불이며, 앞에 \code{<}나 \code{>}를 붙여 바이트 순서를 밝힐 수 있다. 내보내는 쪽은 버퍼가 돌려질 때까지 메모리를 옮기거나 해제하지 않아야 한다. GPU 메모리까지 아우르는 비슷한 약속으로 DLPack도 있다\cite{dlpack}.

가져오는 쪽은 형식이 자기 자료형과 정확히 같을 때만 메모리를 공유할 수 있다. 정수의 폭이 다르거나 바이트 순서가 다르거나 원소가 정렬되어 있지 않으면 변환하면서 복사해야 한다. NEON의 읽기 명령은 정렬되지 않은 주소도 받지만, C++에서 정렬되지 않은 주소를 \code{double*}로 읽는 것은 정의되지 않은 동작이고 컴파일러는 정렬을 가정하고 코드를 만들므로, 정렬되지 않은 메모리는 공유하지 않고 복사하는 것이 안전하다.

\begin{instardust}[\code{docs/design/array.md} §12.3]
stardust의 배열은 버퍼 프로토콜을 직접 구현해 형식 \code{?}, \code{q}, \code{d}, \code{Zd}로 메모리를 내보내므로, numpy와 \code{memoryview}가 복사 없이 읽는다. 가져올 때는 형식이 저장 형태와 정확히 같고 원소가 정렬되어 있으면 공유하고, \code{int32}에서 \code{int64}로처럼 정보를 잃지 않는 확장이면 변환하며 복사하고, big-endian이면 바이트를 뒤집으며 복사한다. 불 버퍼에 $0$과 $1$이 아닌 바이트가 있으면 정규화해서 복사하는데, 모든 바이트의 논리합이 $1$ 이하인지 보는 벡터 리덕션 하나로 검사한다. numpy 배열만은 버퍼 프로토콜을 거치지 않고 numpy의 공개 C 구조체에서 주소와 모양과 보폭을 바로 읽는다. 설계 문서에 따르면 numpy는 내보낼 때마다 형식 문자열을 다시 만들어 호출 시간의 절반을 썼고, 바로 읽는 길이 약 $3$배 빨랐다.
\end{instardust}

\subsection{The Cost of Small Calls}

큰 배열에서는 계산이 시간을 정하지만, 원소가 몇 개뿐인 배열에서는 경계를 넘는 고정 비용이 시간을 정한다. Python에서 확장 함수를 한 번 부르면 인자를 Python 객체에서 C++ 값으로 바꾸고, 결과를 담을 Python 객체를 할당하고, 그 객체를 확장 모듈의 자료형으로 초기화하는 일이 따른다. 이 비용은 원소 하나를 더하는 비용의 수백 배에 이를 수 있다. 작은 배열을 많이 다루는 코드, 예컨대 반복문 안에서 작은 벡터를 계산하는 코드에서는 이 고정 비용을 줄이는 것이 계산을 빠르게 하는 것보다 중요하다.

고정 비용을 줄이는 방법은 대체로 ``할당을 피하고, 찾기를 캐시하고, 필요한 만큼만 만지기''이다. 해제된 결과 객체를 버리지 않고 모아 두었다가 다시 쓰면 할당과 해제가 사라지고, 자료형을 알아보는 검사를 매번 해시표에서 찾는 대신 포인터 하나와 비교하면 찾기가 사라지며, 고정 크기 배열을 통째로 복사하지 않고 실제로 쓰는 칸만 옮기면 메모리 복사가 줄어든다. 작은 데이터를 별도의 할당 없이 객체 안에 담는 것도 같은 종류의 기법이다.

\begin{instardust}[\code{docs/design/array.md} §11]
stardust의 C++/Python 바인딩은 nanobind로 만든다\cite{nanobind}. 설계 문서는 작은 배열의 호출을 프로파일러로 본 결과를 적어 두는데, \code{NDArray}를 담은 Python 객체가 $560$바이트여서 CPython의 작은 객체 할당기를 쓰지 못했고, nanobind가 인스턴스마다 해시표에 넣고 빼는 일과 자료형을 찾는 일이 시간의 상당 부분이었다. 그래서 해제된 배열 객체 $256$개까지를 해시표에 등록된 채로 모아 두었다가 다시 쓰고, 배열 자료형의 포인터를 저장해 두어 \code{Py\_TYPE(o)}와 바로 비교하며, $256$바이트 이하의 데이터는 버퍼 객체 안에 담아 할당을 한 번으로 줄였다. 그 결과 작은 원소별 연산은 numpy보다 약 $2$배 빨라졌다고 설계 문서는 기록한다.
\end{instardust}

\section{Concluding Remarks}

두 언어로 된 라이브러리의 경계는 단순한 연결부가 아니다. C++ 쪽은 템플릿과 RAII로 많은 코드를 안전하게 만들되 정의되지 않은 동작과 별칭 규칙이라는 함정을 피해야 하고, Python 쪽은 참조 계수와 GIL이라는 규칙을 지켜야 하며, 경계 자체는 메모리를 복사 없이 넘기는 약속과 작은 호출의 고정 비용이라는 문제를 안고 있다. 2권에서 다룬 기계의 이야기는 여기서 끝난다. 3부부터는 이 기계 위에서 돌아가는 수치 알고리즘, 먼저 선형대수의 알고리즘으로 넘어간다.

\chapterreferences
